\subsection{Every elliptic Hopf surface occurs unchanged}

\begin{statement}{C.2.UN.37}{Proved}{Occurrence of every Hopf surface of algebraic dimension one}
Let $H$ be any primary or secondary Hopf surface with $a(H)=1$.
There is a proper holomorphic map from a smooth compact complex
threefold to $\mathbb P^1$, with connected fibres and smooth
two-tori away from one value, whose fibre over that value contains
$H$ itself as a smooth component of multiplicity one. Its reduced
special fibre has simple normal crossings and Euler number zero.
Over the complement of the special value the family is the
actual product of two fixed elliptic curves with the base.

The map used to apply C.2.UN.34 is the actual connected elliptic
reduction
\[
                         q:H\longrightarrow\mathbb P^1.
\tag{C.EH.1}
\]
Every fibre of $q$ is a multiple of a smooth elliptic curve,
the modulus is constant, and the integral $H^1$ monodromy is
the identity. The assertion includes secondary Hopf surfaces
without imposing a cyclic condition on their finite primary cover.
\end{statement}

\begin{proof}
We first construct (C.EH.1) and check its full fibres directly
on the given surface. For any Hopf surface one has
\[
                    b_1(H)=1,\qquad H^2(H,\mathbb Q)=0.
\tag{C.EH.2}
\]
These statements follow through the actual finite unramified
primary cover: the primary cover has the rational cohomology
of $S^1\times S^3$, transfer gives $b_1\leq1$ and $b_2=0$,
and the nonzero deck character
$\gamma\mapsto\log|\det(d\gamma)_0|$ gives $b_1\geq1$.
This is the comparison proved in (C.HQ.5)--(C.HQ.7).
Thus it applies to secondary surfaces as well.

Choose a nonconstant meromorphic function $h$ on $H$.
Its effective zero and pole divisors have no common
irreducible component and have intersection number zero
by (C.EH.2). Local intersection multiplicities of distinct
effective curves on a smooth surface are positive.
Consequently those two divisors are disjoint. In a local
coprime expression $h=f/g$, an indeterminacy point would
belong to both $(f=0)$ and $(g=0)$, so no such point exists.
The function therefore defines an actual holomorphic map
$h:H\to\mathbb P^1$. Its Stein factorization is
\[
               H\xrightarrow{q}B\xrightarrow{\nu}\mathbb P^1,
\tag{C.EH.3}
\]
where $B$ is a smooth compact connected curve, $q$ has
connected fibres, and $\nu$ is finite. The degree-one
Leray edge map injects $H^1(B,\mathbb Q)$ into
$H^1(H,\mathbb Q)$, since $q_*\mathbb Q=\mathbb Q$.
Thus $2g(B)\leq1$, and $B\simeq\mathbb P^1$.

Every irreducible compact curve $C\subset H$ is smooth
elliptic. Indeed, adjunction and (C.EH.2) give
\[
                 2p_a(C)-2=(K_H+C)\cdot C=0.
\tag{C.EH.4}
\]
Its normalization cannot be rational: a holomorphic map
$\mathbb P^1\to H$ lifts through the universal cover
$\mathbb C^2\setminus\{0\}\to H$, and both coordinate
functions of the lift are constant. If $g$ is the genus
of the normalization and $\delta$ is its normalization
defect, (C.EH.4) says $g+\delta=1$. Hence $g=1$ and
$\delta=0$. Furthermore two distinct irreducible curves
are disjoint, again by their zero intersection number
and local positivity. Since $q$ has connected fibres,
each full fibre is therefore
\[
                              q^{-1}(b)=m_b C_b,
\tag{C.EH.5}
\]
with $C_b$ a smooth connected elliptic curve. The scheme
fibre is the displayed Cartier divisor: a map from the
smooth surface to a curve has no embedded fibre component.

We check the monodromy by an actual finite map, including
at a multiple fibre. Take a small disc with coordinate $t$
at $b$, write $m=m_b$, and normalize the base change
$t=s^m$. If $x=0$ defines $C_b$ locally, then
$t=u x^m$ for a holomorphic unit $u$. The normalized
cover is given locally by
\[
                     w^m=u,\qquad s=wx.
\tag{C.EH.6}
\]
It is a finite unramified cover of the original surface
tube: locally it is the disjoint union of the $m$ choices
of a holomorphic root of $u$. Denote it by $V\to U$.
The deck group $\mu_m$ acts freely, and $V/\mu_m=U$.
The punctured base-changed family has connected fibres
and connected total space. The normalization has no
component supported over zero. Its Stein factor over
$\Delta_s$ is finite and generically degree one and
therefore is the disc itself. In particular the central
fibre $C'$ is connected. It is smooth by (C.EH.6), and
its multiplicity as the fibre of $V\to\Delta_s$ is one.
Thus this map is a submersion on a neighbourhood of $C'$.
The actual central comparison is the finite unramified
quotient
\[
                         C'\longrightarrow C'/\mu_m=C_b.
\tag{C.EH.7}
\]

A fixed-point-free finite-order automorphism of an elliptic
curve is a translation. To see this, choose an origin and
write it as $z\mapsto\beta z+a$. If $\beta\ne1$, the
nonzero elliptic endomorphism $1-\beta$ is an isogeny,
so $(1-\beta)z=a$ has a solution. Hence every nonidentity
element of the deck group in (C.EH.7) is a translation.
Choose a $\mu_m$-invariant smooth horizontal connection
for the submersion $V\to\Delta_s$, by finite averaging.
Radial parallel transport then gives an equivariant
smooth trivialization. A positive meridian in the
$t$-disc acts on $H^1$ by the corresponding deck
translation or its inverse, according to the transport
convention; both induce the identity. We have proved
\[
                  T_b=I\quad\hbox{on }H^1(q^{-1}(t),\mathbb Z).
\tag{C.EH.8}
\]
The finitely many meridians generate the fundamental
group of the punctured base $\mathbb P^1$, so the entire
integral monodromy representation is trivial.

The holomorphic elliptic $j$-function extends across
each omitted value because the actual cover
$V\to\Delta_s$ is a smooth elliptic family there;
its $j$-function is deck invariant and descends to
the $t$-disc. Thus $j$ is holomorphic on $\mathbb P^1$
with values in $\mathbb C$ and is constant. The map
(C.EH.1) has exactly the fixed-modulus, identity-monodromy,
smooth-multiple-fibre structure required in C.2.UN.34.
Apply that theorem to the given surface $H$ and this
specific map $q$. Its root-curve torsor is the one
defined by (C.EH.6)--(C.EH.7), and its complete analytic
torsor class is retained by the two-line-bundle lift.
Consequently the distinguished component in that
construction is isomorphic to $H$, with multiplicity
one, rather than merely having its topological or
birational invariants. The proper compact map, exact
punctured product and Euler-zero simple-normal-crossing
support are those constructed in (C.SF.8)--(C.SF.13).
\end{proof}
